\documentclass[11pt,a4paper]{article}
\usepackage[utf8]{inputenc}
\usepackage[T1]{fontenc}
\usepackage[ngerman]{babel}
\usepackage{graphicx}

\usepackage{helvet}
\renewcommand{\familydefault}{\sfdefault}

\usepackage{geometry}
\geometry{a4paper, left=2.5cm, right=2.5cm, top=3cm, bottom=3cm}

\usepackage{xcolor}
\usepackage[most]{tcolorbox}
\usepackage{hyperref}
\usepackage{tikz}

\definecolor{uniblue}{RGB}{0, 90, 160}
\definecolor{terminalbg}{RGB}{30, 30, 30}
\definecolor{terminalheader}{RGB}{50, 50, 50}
\definecolor{promptgreen}{RGB}{78, 154, 6}
\definecolor{promptblue}{RGB}{52, 101, 164}
\definecolor{warnred}{RGB}{204, 0, 0}
\definecolor{warnbg}{RGB}{255, 240, 240}

\hypersetup{
	colorlinks=true,
	linkcolor=uniblue,
	urlcolor=uniblue
}

\newcommand{\macbuttons}{%
	\tikz[baseline=-0.6ex]{
		\fill[red!80!white] (0,0) circle (3.5pt);
		\fill[yellow!90!white] (0.4,0) circle (3.5pt);
		\fill[green!80!white] (0.8,0) circle (3.5pt);
	}%
}

\newtcolorbox{terminalbox}[1][]{
	enhanced,
	colback=terminalbg,
	colframe=terminalheader,
	colbacktitle=terminalheader,
	coltext=white,
	fontupper=\ttfamily\small,
	arc=6pt,
	boxrule=1pt,
	toptitle=4pt, bottomtitle=4pt,
	left=12pt, right=12pt, top=10pt, bottom=10pt,
	title={\macbuttons\hspace*{0.5cm}\textcolor{gray!40}{\textsf{\textbf{#1}}}},
	drop shadow={black!30}
}

\newcommand{\prompt}{\textcolor{promptgreen}{student@linux} \textcolor{promptblue}{\char`\~ \$} }

\newtcolorbox{infobox}[1][Wichtiger Hinweis]{
	enhanced,
	colback=warnbg,
	colframe=warnred,
	borderline west={4pt}{0pt}{warnred}, 
	fonttitle=\bfseries\sffamily,
	title={#1},
	coltitle=warnred,
	attach title to upper,
	after title={\\[0.5em]},
	boxrule=0pt,
	arc=3pt,
	left=12pt, right=12pt, top=8pt, bottom=8pt
}

\newcommand{\step}[2]{
	\vspace{0.6cm}
	\noindent\fcolorbox{uniblue}{uniblue}{\textcolor{white}{\textbf{\textsf{ #1 }}}} \textcolor{uniblue}{\textbf{\Large #2}}
	\vspace{0.2cm}\newline
}

\begin{document}
	
	\begin{center}
		{\Huge \textbf{\textsf{\textcolor{uniblue}{CrypTool 2 Web-Umgebung}}}}\\[0.3cm]
		{\large \textsf{Verbindungs- und Nutzungsanleitung}}\\[0.4cm]
		\noindent\textcolor{uniblue}{\rule{\textwidth}{1.5pt}}
	\end{center}
	Diese Anleitung richtet sich an alle, auf deren Geräten die Cryptool-sh Installation nicht funktioniert darunter Chromebooks, ARM-Prozessoren (z. B. macOS).
	
	\step{Schritt 1}{Verbindung zum Server aufbauen}
	In die Adressleiste des Browsers \href{https://132.199.121.49/}{\texttt{https://132.199.121.49/}} eingeben. Sicherheitswarnung müssen eventuell, durch Aufruf des Buttons „Erweiterte Optionen“ umgangen werden. Danach die gelb markierte Schaltfläche bestätigen.
	
	\vspace{0.3cm}
	\begin{center}
		\includegraphics[width=0.8\textwidth]{/home/waznypatrick/Cryptool2/Dokumentation/Bilder/sicherheitswarung_https.png}
	\end{center}
	\vspace{0.3cm}
	
	\begin{infobox}[]
		Falls die Website keine Daten sendet: Überprüfen ob Adresse mit https:// beginnt.
	\end{infobox}
	\step{Schritt 2}{Beim Server anmelden}
	\begin{itemize}
		\item \textbf{Benutzer:} Nach Git hochladen
		\item \textbf{Passwort:} Nach Git hochladen
	\end{itemize}
	\step{Schritt 3}{Eigene Session starten}
	\begin{itemize}
		\item Button in der Mitte \textbf{„Cryptool 2“} auswählen.
		\item Starten und kurz warten.
	\end{itemize}
	
	\noindent\textcolor{uniblue}{\rule{\textwidth}{1.5pt}}
	
	\step{Schritt 4}{Dateien importieren}
	\begin{minipage}[c]{0.55\textwidth}
		\begin{itemize}
			\item Markierten Pfeil anklicken, um das Menü zu öffnen.
			\item Hochladen auswählen und Dateien auswählen.
			\item Danach in Cryptool Öffnen wählen.
			\item Den Ordner Uploads auswählen.
		\end{itemize}
	\end{minipage}
	\hfill
	\begin{minipage}[c]{0.4\textwidth}
		\begin{center}
			\includegraphics[width=\textwidth]{/home/waznypatrick/Cryptool2/Dokumentation/Bilder/importiere_datei.png}
		\end{center}
	\end{minipage}
	\vspace{0.3cm}
	
	\noindent\textcolor{uniblue}{\rule{\textwidth}{1.5pt}}
	
	\step{Schritt 5}{Dateien exportieren}
	\begin{itemize}
		\item In Cryptool Speichern wählen.
		\item Die Datei im Ordner Downloads speichern.
		\item Kasm Menü öffnen.
		\item Markierten Pfeil anklicken, um das Menü zu öffnen.
		\item Hier Herunterladen auswählen.
	\end{itemize}
	
\end{document}